\chapter{Trees and Finite-Difference Pricers in Practice}\label{ch:dv:trees-and-finite-difference-pricers-in-practice}

A pricer returns 6.0864 for a one-year American put with a 200-step binomial tree and 6.0975 with 201 steps. The trader
asks which to trust. The answer is neither: the true value, to five decimals, is 6.0904, below one and above the other. The
\omterm{def:dv:the-binomial-model:crr}{Cox--Ross--Rubinstein tree} of chapter 2 converges to it, but in a zigzag that alternates with the parity of the step count,
and only at first order. At 1\,500 steps both neighbours are finally within a tenth of a cent. A Crank--Nicolson grid with the
right start-up and the strike on a node gets there with 200 nodes and 200 time steps, 28 times less work. This chapter is about
the difference between a pricer that converges and one that is accurate. It covers the pricing equation on a grid, the
conditions that dividends and barriers impose, American exercise by a direct algorithm, and how to spend a computing budget.

\section{The pricing equation on a grid}

For one underlying the \omterm{def:dv:black-scholes-three-ways:pde}{Black--Scholes equation} (chapter 3) is solved most naturally in $x=\ln S$ and in the time to expiry
$\tau$:
\[
\partial_\tau V=\tfrac12\sigma^2\partial_{xx}V+\bigl(r-q-\tfrac12\sigma^2\bigr)\partial_xV-rV.
\]
On a grid $x_0<\dots<x_M$ the derivatives become three-point differences. The build allows the spacing to vary, with
$h_\pm$ the distances to the neighbours, so that nodes can be concentrated near the strike where the payoff has its kink.
In time the $\theta$-scheme (One Quant Book 4, chapter 27) mixes explicit and implicit steps. $\theta=\frac12$ is
Crank--Nicolson, second-order accurate and unconditionally stable, but not free of oscillations. Each step solves one
tridiagonal system, in $O(M)$ operations.

Crank--Nicolson has one known weakness, and option payoffs trigger it. The scheme damps high-frequency components of the
solution only weakly. A kinked payoff contains all frequencies, so the first steps leave oscillations near the strike that
decay slowly. Prices barely notice them, and gammas do. The standard remedy is Rannacher's start-up: a few fully implicit half
steps first, which damp the high frequencies, then Crank--Nicolson. A second remedy works on the payoff itself.

\begin{definition}[Payoff smoothing]\label{def:dv:trees-and-finite-difference-pricers-in-practice:smoothing}
\emph{Payoff smoothing}\index{payoff smoothing} replaces the payoff's value at each grid node by its average over the node's
cell, so that the initial condition carries the kink's effect without its discontinuous derivative; with Rannacher start-up
it restores the scheme's second-order convergence for non-smooth payoffs.
\end{definition}

\begin{example}[Gamma near expiry]\label{ex:dv:trees-and-finite-difference-pricers-in-practice:gamma}
A European put one week from expiry (strike 100, 20\% volatility), on a 200-node grid with ten time steps. Plain
Crank--Nicolson gives a gamma of 0.217 at the strike, where the exact value is 0.141, and wiggles around the true curve on
either side. With two Rannacher start-up steps the largest error is 0.0002
(\cref{fig:dv:trees-and-finite-difference-pricers-in-practice:gamma}). A risk system that reports Crank--Nicolson gammas
without the start-up will show hedgers a phantom gamma spike at every strike near expiry.
\end{example}

\begin{omfigure}
\begin{tikzpicture}
\begin{axis}[width=10.5cm,height=5cm,
  xlabel={spot},ylabel={gamma},
  tick label style={font=\small},label style={font=\small},
  xmin=90,xmax=110,ymin=0,ymax=0.24,grid=major,grid style={gray!25},
  yticklabel style={/pgf/number format/fixed,/pgf/number format/precision=2},
  legend columns=3,legend cell align=left,legend style={font=\footnotesize,at={(0.5,-0.3)},anchor=north,draw=none,fill=none,/tikz/every even column/.append style={column sep=8pt}},
  table/col sep=comma]
\addplot[gray,very thick] table[x=s,y=exact]{figdata/derivatives/22-trees-and-finite-difference-pricers-in-practice/gamma.csv};
\addplot[omProp,thick,mark=*,mark size=1pt] table[x=s,y=cn]{figdata/derivatives/22-trees-and-finite-difference-pricers-in-practice/gamma.csv};
\addplot[omThm,very thick,dashed] table[x=s,y=rannacher]{figdata/derivatives/22-trees-and-finite-difference-pricers-in-practice/gamma.csv};
\legend{exact,Crank--Nicolson,with Rannacher start-up}
\end{axis}
\end{tikzpicture}

\omcaption{Gamma of a one-week put from a 200-node grid and ten time steps. Plain Crank--Nicolson leaves an oscillation at the
strike that ten steps cannot damp; two implicit half steps at the start remove it. Data: the
tutorial.}\label{fig:dv:trees-and-finite-difference-pricers-in-practice:gamma}
\end{omfigure}

Trees are the same equation on a special grid. The binomial tree ties the space step to the time step, $\Delta x=\sigma\sqrt{\Delta t}$,
and is an explicit scheme. The \omterm{def:dv:trees-and-finite-difference-pricers-in-practice:trinomial}{trinomial tree} loosens the tie.

\begin{definition}[Trinomial tree]\label{def:dv:trees-and-finite-difference-pricers-in-practice:trinomial}
A \emph{trinomial tree}\index{trinomial tree} moves the log-price up, down or not at all at each step, by $\Delta x$ (for instance
$\sigma\sqrt{3\Delta t}$), with probabilities chosen to match the mean and variance of the log-return; it is an explicit finite-difference
scheme with a stability condition built into its geometry, and its extra middle branch lets the node grid be placed more freely than a
binomial tree's.
\end{definition}

\section{Discrete dividends and jump conditions}

A cash dividend $D$ paid at $t_d$ makes the share fall by $D$ on the ex-date while the option's value is continuous, since nothing is
paid to its holder. The grid enforces this through a condition between the two sides of the date.

\begin{definition}[Jump condition]\label{def:dv:trees-and-finite-difference-pricers-in-practice:jump}
A \emph{jump condition}\index{jump condition} links a pricing equation's solution across a date at which the state variable jumps
deterministically: for a cash dividend, $V(t_d^-,S)=V(t_d^+,S-D)$, applied on the grid by interpolation, followed by a fresh Rannacher
start-up, since the interpolated profile is no longer smooth where American exercise bites.
\end{definition}

A tree cannot do this directly: shifting a \omterm{def:dv:the-binomial-model:recombining}{recombining tree}'s nodes by a cash amount stops it recombining, which is why chapter 5 used the
escrowed model. The grid has no such problem.

\begin{example}[A three-unit dividend in six months]\label{ex:dv:trees-and-finite-difference-pricers-in-practice:dividend}
With a dividend of 3 paid at six months, the one-year European put is worth 6.835 on the grid, against 6.719 from Black--Scholes on the spot
less the dividend's present value, the escrowed approximation at the same volatility. The American put is worth 7.407, against 6.090
without the dividend.
\end{example}

\section{Barriers and grid alignment}

A continuously monitored barrier is a boundary condition: the option is worth zero, or its rebate, at the barrier. The grid must therefore
have a node exactly at the barrier. If the barrier falls between nodes, the pricer either moves it to the nearest node, which is a
first-order error in the barrier level, or interpolates, which works only approximately. The stretched grid of the build starts at the
barrier and concentrates nodes at the strike. Discretely monitored barriers are applied as a projection at the monitoring dates (chapter
15's shift is the analytical counterpart).

\begin{example}[Aligned and misaligned]\label{ex:dv:trees-and-finite-difference-pricers-in-practice:barrier}
A one-year down-and-out call (strike 100, barrier 90) is worth 8.665 in closed form. With the barrier on a node, the grid's error is 0.0049
with 40 nodes and 0.0002 with 200. With uniform grids on which the barrier falls between nodes, the error is 1.10 with 40 nodes, 0.16 with 200 and
0.046 with 300 (\cref{fig:dv:trees-and-finite-difference-pricers-in-practice:barrier}): hundreds of times worse at the same cost, and
dependent on where the barrier happens to fall between two nodes.
\end{example}

\begin{omfigure}
\begin{tikzpicture}
\begin{axis}[width=10.5cm,height=5cm,xmode=log,ymode=log,
  xlabel={grid nodes},ylabel={absolute error},
  tick label style={font=\small},label style={font=\small},
  xmin=35,xmax=340,ymin=5e-5,ymax=3,grid=major,grid style={gray!25},
  xtick={40,60,100,200,300},xticklabels={40,60,100,200,300},
  legend columns=2,legend cell align=left,legend style={font=\footnotesize,at={(0.5,-0.3)},anchor=north,draw=none,fill=none,/tikz/every even column/.append style={column sep=8pt}},
  table/col sep=comma]
\addplot[omThm,very thick,mark=*,mark size=1.6pt] table[x=m,y=aligned]{figdata/derivatives/22-trees-and-finite-difference-pricers-in-practice/barrier.csv};
\addplot[omProp,very thick,mark=triangle*,mark size=1.8pt] table[x=m,y=misaligned]{figdata/derivatives/22-trees-and-finite-difference-pricers-in-practice/barrier.csv};
\legend{barrier on a node,barrier between nodes}
\end{axis}
\end{tikzpicture}

\omcaption{Error of a one-year down-and-out call (strike 100, barrier 90) on grids of increasing size, with the barrier exactly on a node and
between nodes. Alignment, not size, decides the accuracy. Data: the tutorial.}\label{fig:dv:trees-and-finite-difference-pricers-in-practice:barrier}
\end{omfigure}

\section{American exercise}

Early exercise turns the pricing equation into a free-boundary problem (chapter 6): $V\ge g$, the exercise value, everywhere, with the
equation holding where $V>g$. On a grid, the naive treatment solves the linear system and then takes the maximum with $g$. That works with an
implicit step, and it costs Crank--Nicolson its accuracy, because the projection and the solve are not consistent. Iterative methods such as
projected successive over-relaxation solve the constrained system properly, at many times the cost. For a put, whose exercise region is a
single interval of low spots, there is a direct method.

\begin{definition}[Brennan--Schwartz algorithm]\label{def:dv:trees-and-finite-difference-pricers-in-practice:bs}
The \emph{Brennan--Schwartz algorithm}\index{Brennan--Schwartz algorithm} solves a tridiagonal system with the constraint $V\ge g$ exactly
when the exercise region is an interval at one end of the grid: it eliminates from the other end, then back-substitutes towards the exercise
region, taking the maximum with $g$ at each node; its cost is that of one tridiagonal solve.
\end{definition}

The order of elimination is the whole trick. For a put the exercise region lies at low spots, so the elimination runs from the top of the grid
down, and the back substitution runs upwards from the exercise region. There each value is fixed by the constraint before the continuation
values above it are computed. For a call on a dividend-paying share the order reverses.

The grid also gives the \omterm{def:dv:american-options-and-early-exercise:boundary}{exercise boundary} as a by-product: the highest spot at which the value equals the exercise value, at each date,
refined between nodes. The tree gives it too, but only as a staircase of its own nodes
(\cref{fig:dv:trees-and-finite-difference-pricers-in-practice:boundary}). For the one-year put the grid places the boundary at 80.85 one year
before expiry, while a 2\,000-step tree's boundary jumps between 80.68 and 81.04 from one step to the next.

\begin{omfigure}
\begin{tikzpicture}
\begin{axis}[width=10.5cm,height=5cm,
  xlabel={time to expiry (years)},ylabel={exercise boundary},
  tick label style={font=\small},label style={font=\small},
  xmin=0,xmax=1,ymin=78,ymax=101,grid=major,grid style={gray!25},
  legend columns=2,legend cell align=left,legend style={font=\footnotesize,at={(0.5,-0.3)},anchor=north,draw=none,fill=none,/tikz/every even column/.append style={column sep=8pt}},
  table/col sep=comma]
\addplot[omProp,thin] table[x=tau,y=s]{figdata/derivatives/22-trees-and-finite-difference-pricers-in-practice/boundary_tree.csv};
\addplot[omThm,very thick,mark=*,mark size=1.2pt] table[x=tau,y=s]{figdata/derivatives/22-trees-and-finite-difference-pricers-in-practice/boundary_grid.csv};
\legend{2\,000-step binomial tree,grid (400 nodes)}
\end{axis}
\end{tikzpicture}

\omcaption{The American put's \omterm{def:dv:american-options-and-early-exercise:boundary}{exercise boundary} (strike 100, 5\%, 20\%) by time to expiry: exercise below the curve. The tree's boundary is a
staircase of its nodes; the grid's, refined between nodes, is smooth. Data: the tutorial.}\label{fig:dv:trees-and-finite-difference-pricers-in-practice:boundary}
\end{omfigure}

\section{Accuracy against speed}

Every pricer has two dials, the grid and the time steps, and a budget. The question the trader asked in the introduction is really
this: how much computation buys a given accuracy?

\begin{example}[A tenth of a cent]\label{ex:dv:trees-and-finite-difference-pricers-in-practice:budget}
The one-year American put (spot and strike 100, rate 5\%, volatility 20\%) is worth 6.0904, from a Richardson-extrapolated grid (1\,600 and 800 nodes
and steps). The \omterm{def:dv:the-binomial-model:crr}{Cox--Ross--Rubinstein tree}'s error alternates in sign with the parity of the step count. Its two neighbours are both within 0.1 cent
from 1\,500 steps, 1\,125\,750 node updates. The \omterm{def:dv:trees-and-finite-difference-pricers-in-practice:trinomial}{trinomial tree} converges without the zigzag but at first order: its error is 0.11 cent at 1\,600 steps
and 0.054 at 3\,200. The grid, with a node on the strike, \omterm{def:dv:trees-and-finite-difference-pricers-in-practice:smoothing}{payoff smoothing} and Rannacher start-up, errs by 1.04 cent with 50 nodes and steps, 0.30 with
100, 0.091 with 200 and 0.029 with 400: second order. It is within 0.1 cent at 200, 40\,600 node updates. That is 28 times less work than the tree
(\cref{fig:dv:trees-and-finite-difference-pricers-in-practice:work}). Without smoothing and start-up the grid errs by 0.125 cent at 200.
\end{example}

\begin{omfigure}
\begin{tikzpicture}
\begin{groupplot}[group style={group size=2 by 1,horizontal sep=1.6cm},
  width=6.6cm,height=5cm,
  tick label style={font=\small},label style={font=\small},grid=major,grid style={gray!25},table/col sep=comma]
\nextgroupplot[xlabel={tree steps},ylabel={American put},xmin=100,xmax=400,ymin=6.075,ymax=6.105,
  yticklabel style={/pgf/number format/fixed,/pgf/number format/precision=3},
  title={\small the binomial zigzag},
  legend columns=1,legend cell align=left,legend style={font=\footnotesize,at={(0.5,-0.32)},anchor=north,draw=none,fill=none}]
\addplot[omProp,thin] table[x=n,y=crr]{figdata/derivatives/22-trees-and-finite-difference-pricers-in-practice/oscillation.csv};
\addplot[black,thick,dashed] table[x=n,y=ref]{figdata/derivatives/22-trees-and-finite-difference-pricers-in-practice/oscillation.csv};
\legend{Cox--Ross--Rubinstein,value}
\nextgroupplot[xmode=log,ymode=log,xlabel={work (node updates)},ylabel={error (cents)},xmin=500,xmax=2e7,ymin=0.01,ymax=10,
  title={\small error against work},
  legend columns=1,legend cell align=left,legend style={font=\footnotesize,at={(0.5,-0.32)},anchor=north,draw=none,fill=none}]
\addplot[omProp,thick,mark=*,mark size=1pt] table[x=work,y=err]{figdata/derivatives/22-trees-and-finite-difference-pricers-in-practice/work_crr.csv};
\addplot[omDef,thick,mark=square*,mark size=1.4pt] table[x=work,y=err]{figdata/derivatives/22-trees-and-finite-difference-pricers-in-practice/work_tri.csv};
\addplot[omThm,very thick,mark=triangle*,mark size=1.8pt] table[x=work,y=err]{figdata/derivatives/22-trees-and-finite-difference-pricers-in-practice/work_fd.csv};
\addplot[black,dashed] coordinates {(500,0.1) (2e7,0.1)};
\legend{binomial (worst of $n$ and $n+1$),trinomial,grid (smoothing and Rannacher),0.1 cent}
\end{groupplot}
\end{tikzpicture}

\omcaption{A one-year American put. Left: the \omterm{def:dv:the-binomial-model:crr}{Cox--Ross--Rubinstein tree} against the value, by number of steps. Right: absolute error in cents
against work, in node updates, for the binomial tree (the larger error of each pair $n$, $n+1$), the \omterm{def:dv:trees-and-finite-difference-pricers-in-practice:trinomial}{trinomial tree} and the Crank--Nicolson
grid. Data: the tutorial.}\label{fig:dv:trees-and-finite-difference-pricers-in-practice:work}
\end{omfigure}

Work in node updates is a fair proxy for run time between methods written in the same language, and the grid's per-node work (a tridiagonal
solve) is a few times the tree's (a weighted average), so the real speed-up is smaller than 28, though of the same order. Two further economies
follow from second-order convergence. Richardson extrapolation (Book 4) of two grid sizes cancels the leading error term. Grid stretching puts the
nodes where the payoff has its kink. The build carries both. Trees keep a place where their simplicity matters more than their cost: teaching,
checking, and products whose early-exercise rules are easiest to state node by node.

\section{Tutorial: a finite-difference engine}\label{tut:dv:trees-and-finite-difference-pricers-in-practice:tut}

\begin{tutorial}
\textbf{Goal.} Build a Crank--Nicolson pricer on a stretched log-spot grid with Rannacher start-up, \omterm{def:dv:trees-and-finite-difference-pricers-in-practice:smoothing}{payoff smoothing}, dividend \omterm{def:dv:trees-and-finite-difference-pricers-in-practice:jump}{jump conditions},
barrier boundaries and American exercise by Brennan--Schwartz; compare it with the trees on error against work. \textbf{End state:} the four
figures and the numbers of the weekend problem.
\begin{tutsteps}
\item \textbf{The operator on a non-uniform grid}: three-point differences with unequal spacings:
\omcode{code/firm/fdpricer/firm_fdpricer.py}{42}{50}{Tridiagonal coefficients of the pricing operator.}\label{lst:dv:trees-and-finite-difference-pricers-in-practice:operator}
\item \textbf{Brennan--Schwartz}: eliminate from the top, back-substitute upwards with the constraint:
\omcode{code/firm/fdpricer/firm_fdpricer.py}{68}{82}{The Brennan--Schwartz solver.}\label{lst:dv:trees-and-finite-difference-pricers-in-practice:bs}
\item \textbf{The time loop}: Rannacher half steps, Crank--Nicolson, boundaries, dividends:
\omcode{code/firm/fdpricer/firm_fdpricer.py}{107}{135}{The theta-scheme time loop with jump conditions.}\label{lst:dv:trees-and-finite-difference-pricers-in-practice:loop}
\item \textbf{The C++20 twin} of the solver, tested against the Python numbers to rounding (the Rust twin has the same algorithm):
\omcode{code/firm/fdpricer/cpp/firm_fdpricer.hpp}{36}{52}{Brennan--Schwartz in C++20.}\label{lst:dv:trees-and-finite-difference-pricers-in-practice:cpp}
\item \textbf{Run} \texttt{dv\_fd.accuracy\_study()}, \texttt{gamma\_profile()}, \texttt{dividend\_example()}, \texttt{barrier\_alignment()} and
\texttt{fig\_fd.py}.
\end{tutsteps}
\textbf{What to change next.} Replace Brennan--Schwartz by solve-then-project with Crank--Nicolson and measure the order of convergence; price the
American call with the dividend and find the grid's \omterm{def:dv:american-options-and-early-exercise:boundary}{exercise boundary} just before the ex-date; time the Python, C++ and Rust versions on the same
grid.
\end{tutorial}

\section{Build: the finite-difference engine}\label{bld:dv:trees-and-finite-difference-pricers-in-practice:fdpricer}

\begin{build}
\bfield{Purpose} The miniature firm's one-dimensional PDE engine: the finite-difference engine of chapter 28's pricing library, used for American
exercise, barriers, dividends and convertibles (chapter 21); twins in C++20 and Rust for the production path.
\bfield{Interface} \texttt{uniform\_grid(s0, vol, t, m, width)}; \texttt{stretched\_grid(s0, k, vol, t, m, width, concentration, lower)};
\texttt{solve(payoff, x, t, r, q, vol, n\_t, american, rannacher, dividends, smoothing, lower, upper) -> values}; \texttt{value\_at(x, v, s0) -> value,
delta, gamma}; \texttt{american\_put}; \texttt{smoothed\_payoff}; \texttt{trinomial(s0, k, t, r, q, vol, n, right, american)}. C++20
\texttt{firm::fdpricer::american\_put}, Rust \texttt{firm\_fdpricer::american\_put}.
\bfield{Rules} The strike (and any barrier) is a grid node; Rannacher start-up after expiry and after every discontinuity; American exercise by
Brennan--Schwartz where the exercise region is an interval, projected iteration otherwise; every price comes with its convergence check (two grid
sizes).
\bfield{Acceptance tests} \texttt{code/firm/fdpricer/tests/}: European value and delta against Black--Scholes on both grids; the American put converges at
second order; the \omterm{def:dv:trees-and-finite-difference-pricers-in-practice:trinomial}{trinomial tree} converges; the dividend \omterm{def:dv:trees-and-finite-difference-pricers-in-practice:jump}{jump condition} against a simulation; a barrier on a node against the closed form;
\texttt{cpp/} and \texttt{rust/}: the twins reproduce the Python engine's numbers to $10^{-9}$ and converge.
\bfield{Stretch} Two-dimensional grids (alternating directions) for stochastic volatility; a projected successive over-relaxation fallback; adaptive time
steps around dividends.
\end{build}

\begin{omsources}
\item M.~J.~Brennan and E.~S.~Schwartz, ``The valuation of American put options'', \emph{Journal of Finance} 32(2) (1977) 449--462.
\item R.~Rannacher, ``Finite element solution of diffusion problems with irregular data'', \emph{Numerische Mathematik} 43(2) (1984) 309--327.
\item D.~M.~Pooley, K.~R.~Vetzal and P.~A.~Forsyth, ``Convergence remedies for non-smooth payoffs in option pricing'', \emph{Journal of
Computational Finance} 6(4) (2003) 25--40.
\item D.~Tavella and C.~Randall, \emph{Pricing Financial Instruments: The Finite Difference Method} (Wiley).
\end{omsources}

\section{Exercises}

\begin{exercise}[$\star$]\label{exo:dv:trees-and-finite-difference-pricers-in-practice:1}
Why does the binomial tree's price alternate above and below the value as the number of steps changes parity?
\end{exercise}

\begin{exercise}[$\star$]\label{exo:dv:trees-and-finite-difference-pricers-in-practice:2}
A grid's error is 0.30 cent with 100 nodes and 0.091 with 200. Estimate the order of convergence and the error with 400.
\end{exercise}

\begin{exercise}[$\star$]\label{exo:dv:trees-and-finite-difference-pricers-in-practice:3}
Why can a \omterm{def:dv:the-binomial-model:recombining}{recombining tree} not handle a cash dividend directly, and how does the grid?
\end{exercise}

\begin{exercise}[$\star\star$]\label{exo:dv:trees-and-finite-difference-pricers-in-practice:4}
Using the grid values with 800 and 1\,600 nodes and steps, derive the Richardson-extrapolated value for a second-order method.
\end{exercise}

\begin{exercise}[$\star\star$]\label{exo:dv:trees-and-finite-difference-pricers-in-practice:5}
Explain why Brennan--Schwartz must eliminate from the top of the grid for a put.
\end{exercise}

\begin{exercise}[$\star\star$]\label{exo:dv:trees-and-finite-difference-pricers-in-practice:6}
Why does moving a barrier by half a grid cell produce a first-order error in the price?
\end{exercise}

\begin{exercise}[$\star\star\star$]\label{exo:dv:trees-and-finite-difference-pricers-in-practice:7}
\emph{Coding.} Price the American put by solving with Crank--Nicolson and projecting afterwards instead of Brennan--Schwartz. Measure the error at 100, 200
and 400 nodes and the order.
\end{exercise}

\begin{exercise}[$\star\star\star$]\label{exo:dv:trees-and-finite-difference-pricers-in-practice:8}
\emph{Find the flaw.} ``Our tree prices the American put to six decimals: we ran it at 2\,000 steps and the result did not change in the sixth decimal
between runs.''
\end{exercise}

\section{Problem: Neither 200 Nor 201}

\begin{problem}[{Weekend problem --- how much computation a tenth of a cent costs}]\label{pb:dv:trees-and-finite-difference-pricers-in-practice:1}
A risk system prices a book of one-year American puts (spot and strike 100, rate 5\%, volatility 20\%) with a 200-step binomial tree. The desk must certify every
price to a tenth of a cent.

\medskip\noindent\textbf{Part I --- The tree.}
\begin{enumerate}
\item Give the tree's price with 200 and 201 steps, and the value.
\item Give the two errors in cents.
\item From how many steps are both neighbours within 0.1 cent?
\item How many node updates is that?
\item Why is averaging $n$ and $n+1$ not a real fix?
\end{enumerate}

\medskip\noindent\textbf{Part II --- The grid.}
\begin{enumerate}[resume]
\item Give the grid's error with 50, 100, 200 and 400 nodes and steps.
\item Which size first meets the tolerance, and at what work?
\item What does the grid lose without \omterm{def:dv:trees-and-finite-difference-pricers-in-practice:smoothing}{payoff smoothing} and Rannacher start-up?
\item Give the \omterm{def:dv:trees-and-finite-difference-pricers-in-practice:trinomial}{trinomial tree}'s errors at 1\,600 and 3\,200 steps.
\item What would Richardson extrapolation of 100 and 200 nodes give?
\end{enumerate}

\medskip\noindent\textbf{Part III --- The rest of the book.}
\begin{enumerate}[resume]
\item The book also holds puts with a dividend of 3 in six months. Give the European and American prices.
\item And down-and-out calls with a barrier at 90: what goes wrong on a uniform grid?
\item Which Greek does plain Crank--Nicolson get wrong near expiry, and by how much in the chapter's example?
\item Why is work in node updates only a proxy for run time?
\item How would the C++ twin change the picture?
\end{enumerate}

\medskip\noindent\textbf{Part IV --- Judgement.}
\begin{enumerate}[resume]
\item What would you tell the trader who asked which of 200 and 201 to trust?
\item How would you certify a price's accuracy in production?
\item When is a tree still the right tool?
\item State the \emph{named result}: the grid size and work at which the American put is within 0.1 cent, for the tree and for the grid, and the speed-up of
the grid.
\item In one sentence: what distinguishes a pricer that converges from one that is accurate?
\end{enumerate}
\end{problem}

\section{Interview questions}

\begin{interviewq}[$\star$ \iqroles{developer, researcher}]\label{iq:dv:trees-and-finite-difference-pricers-in-practice:1}
Why does Crank--Nicolson produce oscillating \omterm{def:dv:greeks-and-the-hedging-pnl:greeks}{Greeks} for options near expiry, and how do you fix it?
\end{interviewq}

\begin{interviewq}[$\star\star$ \iqroles{developer}]\label{iq:dv:trees-and-finite-difference-pricers-in-practice:2}
How do you price an American put on a finite-difference grid efficiently?
\end{interviewq}

\begin{interviewq}[$\star\star$ \iqroles{researcher}]\label{iq:dv:trees-and-finite-difference-pricers-in-practice:3}
Compare the convergence of a binomial tree, a \omterm{def:dv:trees-and-finite-difference-pricers-in-practice:trinomial}{trinomial tree} and a Crank--Nicolson grid for an American put.
\end{interviewq}

\begin{interviewq}[$\star\star$ \iqroles{developer}]\label{iq:dv:trees-and-finite-difference-pricers-in-practice:4}
How do you handle a discrete cash dividend and a continuously monitored barrier on a grid?
\end{interviewq}

\begin{interviewq}[$\star\star$ \iqroles{risk}]\label{iq:dv:trees-and-finite-difference-pricers-in-practice:5}
How would you validate the numerical accuracy of a production pricer?
\end{interviewq}

\begin{interviewq}[$\star\star\star$ \iqroles{developer, researcher}]\label{iq:dv:trees-and-finite-difference-pricers-in-practice:6}
Your grid pricer for an American call on a dividend-paying share gives a slightly wrong \omterm{def:dv:american-options-and-early-exercise:boundary}{exercise boundary} just before the ex-date. Diagnose.
\end{interviewq}
